\section{Discussion\label{sec:discussion}}

The research questions examine where architectural signal resides and how a mapping model can use it. RQ1 compares what file representations encode before message passing, RQ2 adds file, member, and folder context and compares the folder gains with a directory-based rule, and RQ3 examines how the predictor and the direction in which it aggregates dependency neighbors change what they contribute. This section considers the three together and states the conditions under which they hold.

\parg{Dependency information in file representations.}
RQ1~\ref{sec:rq1} compares file representations before the predictor exchanges information between linked files, so the results reflect the architectural information encoded in each representation. Node2Vec, which is learned from the dependency graph alone, without labels, has the highest median micro-F1 and macro-F1 on six of the ten systems. Fine-tuned UniXcoder significantly improves over its zero-shot counterpart on micro-F1 on all ten systems and on macro-F1 on nine. Since the zero-shot encoder receives no dependency signal, this shows that fine-tuning on dependency-derived pairs adds information beyond that of pretrained code representations alone. Dependency information can therefore provide useful architectural signal before a predictor combines information across dependency links, either by encoding graph position directly in Node2Vec or by adapting a contextual code representation through dependency-guided fine-tuning.

The usefulness of this signal varies with the relation between dependency structure and the reference architecture. Node2Vec's micro-F1 is lower on the three least homophilic systems, SH-3D, HDF, and T.Mates (Table~\ref{tab:datasets}), with a median of $58$ compared with $72$ on the remaining seven. When dependencies frequently cross module boundaries, random walks may place files from different modules in similar structural neighborhoods, limiting the module-discriminating information available to Node2Vec. Homophily is computed from the reference mapping and is therefore unavailable before a mapping exists; it is used here to interpret the observed pattern, not to select a representation in practice.

Fine-tuned UniXcoder shows a different pattern on the low-homophily systems. Its median micro-F1 is $67$ on SH-3D, HDF, and T.Mates, compared with $69.5$ on the remaining seven. It exceeds Node2Vec by $9$ and $10$ micro-F1 points on SH-3D and T.Mates, which combine low homophily with the highest average degrees, $\langle k \rangle$ of $23.43$ and $22.90$. In these systems, dependency neighborhoods may span several modules, while the code content represented by UniXcoder can provide complementary evidence within such mixed neighborhoods. HDF is the exception: Node2Vec remains ahead of \ft{} ($57$ versus $48$ micro-F1), although every representation performs comparatively poorly on this system.

The two representations without dependency information, Word2Vec and zero-shot UniXcoder, show no consistent ordering. Their weaker overall results reinforce that the advantage of Node2Vec and \ft{} is associated with dependency information rather than with the use of a graph embedding or a code encoder alone.

The RQ2 and RQ3 comparison concerns a different use of dependencies. While Node2Vec learns file vectors from unlabeled random walks, the neighbor-based predictors use dependency links while learning the mapping from the available labels. The results show that useful dependency information in an input representation does not by itself determine whether this additional neighbor use helps: with N2V, directed aggregation lowers the median macro-F1 for all three predictors, whereas undirected aggregation raises it, and HGT gains least from neighbors with every representation. The effect therefore depends on how the predictor aggregates the neighborhood and in which direction.

\parg{Folder organization and architectural exceptions.}
The gains from adding folder nodes in RQ2 show that directory organization helps reproduce the reference mapping. However, extending a directory assignment is different from identifying files that depart from it. The directory analysis separates these cases by checking, on the files that a directory-based baseline gets wrong, which incorrect predictions of FM the folder nodes correct and which correct predictions they lose (Section~\ref{sec:directory_errors}).

Lucene and HDF mainly illustrate the first case. Their directories are pure under the reference mapping, and directory-baseline errors occur where the immediate directory has no labeled examples. Correcting these errors shows that a model can recover an assignment when local supervision is missing. It gives little evidence about conflicts between directory structure and architecture. Bash has mixed directories and therefore presents a different problem. JabRef shows why the distinction matters: adding folder context improves overall accuracy while losing more correct predictions on directory-baseline errors than it gains.

A mapping tool intended to identify architectural exceptions should therefore be evaluated on files that depart from their directory assignment, alongside overall accuracy. A rule-based directory baseline using only the initial labels provides a useful reference, but the counts compare FM with FFM on the baseline's errors and do not show how each model compares with the baseline on those files. Recovering assignments in directories without initial labels is useful, but does not establish that a model can identify files whose module assignment conflicts with their directory.

\parg{Dependency neighborhood processing and aggregation direction.}
RQ3 distinguishes dependency information encoded in file representations from its additional use during prediction. With undirected aggregation, dependency neighbors raise the median score for GAT, HGT, and the neighbor \mlp{} with every representation, including N2V and \ft{}, so dependency-informed initialization does not make neighbor access redundant. The gains are larger for W2V and \zs{}, which do not encode dependencies, than for N2V and \ft{}, which reinforces the RQ1 finding that dependencies provide signal beyond file content. HGT gains least with every representation and GAT scores highest in most comparisons, although the differences between predictors are small relative to the variation across systems. GAT and HGT differ in their attention formulation and parameterization, so the comparison does not isolate why GAT performs better. These findings concern the evaluated file-level predictors, not the usefulness of dependency types in general or the folder- and member-aware configurations.

The preferred direction of aggregation follows the system and not the predictor or the representation. SH-3D and T.Mates favor directed aggregation, and Lucene, Bash, Chrome, HDC, and HDF favor undirected aggregation, on both metrics. Undirected aggregation also adds the reverse edges, so the comparison does not isolate the direction of propagation from the larger neighborhood. With N2V, directed aggregation lowers the median macro-F1 for all three predictors. The N2V vectors were learned on the undirected dependency graph, which may make directed aggregation less consistent with the neighborhood that each vector encodes, but we did not test this.

\parg{Connectivity and neighbor information.}
Connectivity describes how much neighboring context is available, while homophily describes whether linked files share a module. Connectivity matters particularly for shallow message passing, whose reach is limited by the number of layers. Sparse neighborhoods restrict the features available for aggregation, whereas Node2Vec can encode more distant context through random walks. Connectivity can therefore constrain additional message passing even when the initial structural representation remains useful.

With N2V and undirected aggregation, no predictor changes significantly on Bash and HDC, the systems with the lowest average degree. This is consistent with limited local context reducing the benefit of neighbor aggregation, but does not establish connectivity as the cause. In contrast, T.Mates and SH-3D have low homophily but among the highest average degrees, and they are the systems with the largest N2V micro-F1 gains from file-level message passing in RQ2 ($+7.2$ and $+10.3$ points) and the systems that favor directed aggregation in RQ3. HDF has similarly low homophily but a much lower degree and no comparable gain. These observations suggest that connectivity is relevant to the context available for shallow message passing, but with ten systems and correlated descriptors we cannot separate connectivity from homophily, system size, or label scarcity, and we did not test this.

\parg{Implications for evaluating mapping models.}
Evaluation should establish what the input representation already provides, then isolate neighbor access within a fixed predictor. Comparing a graph model with a graph-free MLP changes both the predictor and its access to structure; switching neighbors on and off within the same predictor tests whether that access helps.

Macro and micro F1 scores should be read together, since correctly mapping more files does not necessarily improve performance across modules. The directory analysis adds a further requirement: aggregate scores should be accompanied by evidence about files that depart from directory-based assignments. Together, these comparisons distinguish the contributions of input features, neighbor processing, and neighbor direction, while exposing limitations hidden by aggregate scores. These points concern the ten systems and the compared configurations and do not rank the signals or the predictors in general.